% 本文件由 LaTeX 排版 Agent 根据有效 translation.json 自动生成。
% author=Justin Martyr; work_id=0130; title=论天主的唯一统治

\chapter{论天主的唯一统治}

% structure: p0001 source=h1 target=omitted reason=page-title-already-rendered-by-parent

% structure: p0002 source=h2 target=section reason=relative-heading-rank; base=h2

\section{第一章　作者之目的}

虽然人性最初被赋予理智与安全的结合，以辨别真理及对万有唯一主当有的敬拜，但嫉妒却暗示人自身伟大的优越性，将人转向制造偶像；这种迷信习俗延续了很长时期后，被传给大多数人，仿佛它是自然且真实的。一个爱人者，或更确切地说，一个爱天主者，其职责是提醒那些忽略此事的人，让他们知道所当知道的事。因为真理本身足以藉着穹苍之下所包含的事物，显明创造它们的那一位所设立的秩序。然而，因着天主的长期容忍，遗忘占据了人的心思，鲁莽地将唯独归于唯一真天主的名称转嫁于凡人；罪的感染从少数人蔓延到多数人，他们被普遍的习俗蒙蔽，以至于不认识那持久不变者。因为前代的人为更有权势者设立了私人和公共的礼仪，使遗忘大公信仰占据了他们的后代；但我，正如我刚才所说，怀着一颗爱天主的心，将采用一位爱人者的言语，并把它放在那些有智慧的人面前——凡有幸观察宇宙管理的人都应当具有这种智慧——以便他们恒久不变地敬拜那知晓万有的天主。我这样做，不是仅靠言辞的炫耀，而是完全使用从希腊文学的古诗中和在所有人中非常普遍的著作中引来的论证。因为藉此，那些将偶像崇拜作为法律传给大众的名人，将被他们自己的诗人和文学作品教导并定明极大的无知。\par

% structure: p0004 source=h2 target=section reason=relative-heading-rank; base=h2

\section{第二章　天主独一的证据}

首先，\Person{埃斯库罗斯}在阐述其作品的结构时，也曾这样论及唯一的天主：\par

\Quote{凡俗之人，请远离圣洁的天主，\newline{}切勿以为祂像你一样，\newline{}披着血肉之衣；因为伟大之神\newline{}对你这样的蠕虫全然未知。\newline{}祂显多样形像；有时\newline{}如同永不满足的焚烧烈火；\newline{}时而如水，时而又\newline{}披上黑暗的衣袍。\newline{}甚至地上的野兽也反映\newline{}祂神圣的形像；而风、云、雨、\newline{}雷声和闪电，\newline{}向人启示他们伟大至高的主。\newline{}在祂面前，大海、岩石、每道泉源\newline{}和所有洪流，都恭敬俯伏；\newline{}当他们凝视祂可畏的容颜，\newline{}高山、大地、海洋的深渊\newline{}和最高的山峰，\newline{}都战栗：因为祂是全能的主；\newline{}这就是至高天主的荣耀。}\par

但不仅他一人领受了关于天主的认识；\Person{索福克勒斯}也这样描述万有唯一创造者、唯一天主的本性：\par

\Quote{有一位天主，的确只有一位，\newline{}祂创造了天和下方广阔的大地、\newline{}海洋的闪光波浪和风；\newline{}但我们许多凡人心中犯错，\newline{}为排解我们的苦难，\newline{}立了石制、铜制的神像，\newline{}或用金、象牙雕刻的像；\newline{}并为这些我们手造的物品\newline{}供奉祭品和盛大的仪式，\newline{}以为这样就是做了一件虔诚的事。}\par

还有\Person{菲勒蒙}，他发表了许多对古代习俗的解释，也分享了对真理的认识；他这样写道：\par

\Quote{告诉我，我们应如何设想天主？\newline{}独一者，看见万有，却自身不可见。}\par

甚至\Person{俄耳甫斯}，虽然他介绍了三百六十位神，也会从名为\WorkTitle{遗嘱}的短篇作品中为我作证，他在那里似乎为他的错误而懊悔，写了以下内容：\par

\Quote{我要对那些合法聆听的人说话；\newline{}其余的人，你们这些亵渎者，请关闭门户！\newline{}啊，\Person{穆赛俄斯}，请听我说，\newline{}你是光明之月的后裔。\newline{}我现在告诉你的话确实是真理，\newline{}如果你曾见过我以往的想法，\newline{}不要让他们夺去你蒙福的生命；\newline{}宁可回转你内心的深处\newline{}到那光明与知识的居所。\newline{}接受圣言引导你的脚步；\newline{}行在笔直确定的道路上，\newline{}仰望那独一而普世的君王，\newline{}独一的、自生的、唯一的那位，\newline{}万有和我们自身都出自于祂。\newline{}万有在祂锐利的凝视下敞开，\newline{}而祂自身依然不可见；\newline{}在祂所有工作中显现，却仍不被看见，\newline{}祂从善中赐予凡人恶，\newline{}降下寒冷战乱和泪流满面的忧伤；\newline{}除那大君王之外别无他者。\newline{}云彩永远环绕祂的宝座；\newline{}凡人的眼目，在凡人的眼中，\newline{}无力看见治理万有的宙斯。\newline{}祂安坐在铜制的天上\newline{}祂的宝座上；祂的脚下\newline{}踏着大地，祂的右手伸展\newline{}到海洋的尽头，四周\newline{}山脉和溪流战栗，\newline{}还有苍茫大海的深处。}\par

他说这些话，仿佛他曾亲眼目睹天主的伟大。\Person{毕达哥拉斯}同意他，写道：\par

\Quote{若有人胆大地说：看哪，我是神！\newline{}除那独一的、永恒的、无限的之外，\newline{}就让他从他篡夺的宝座上\newline{}施展能力，形成另一个世界，\newline{}如同我们所居住的，说：这是我的。\newline{}不仅如此，还要永远居住在这新领域。\newline{}如果他能做到，\newline{}那么他确实可以被宣告为神。}\par

% structure: p0015 source=h2 target=section reason=relative-heading-rank; base=h2

\section{第三章　未来审判的证据}

此外，关于祂，唯独祂有能力对今生所行的作为以及对（人所显示的）对天主的无知施行审判，我可以从你们自己的行列中提出证人；首先是\Person{索福克勒斯}，他这样说：\par

\Quote{那时刻中的时刻将要来到，必将来临，\newline{}当从金色苍穹降下\newline{}火之丰饶宝藏，在烈焰中\newline{}天地万物尽被焚烧；\newline{}然后，当一切受造物溶解，\newline{}海之最后一波将死于岸边，\newline{}光秃大地无树木，燃烧之气\newline{}不再有羽翼之物栖于其怀。\newline{}通往冥府之路有两条，我们深知；\newline{}义者走此，恶者走彼，\newline{}各自趋向分离的命运；而那一位\newline{}曾毁灭万有的，将拯救万有。}\par

\Person{斐利蒙}又说：—\par

\Quote{\Person{尼科斯特拉图}，你以为那些在此\newline{}享受了人生诸多美好之死者，\newline{}能逃脱神明的关注，\newline{}仿佛祂会遗忘他们么？\newline{}不，有一位正义之眼监察万有；\newline{}因为若善人与恶人结局相同，\newline{}那你就去抢、去偷、去掠，随你所欲，\newline{}尽行你眼中看为恶的事。\newline{}但切莫受骗；因为在其下\newline{}已设宝座与审判之处，\newline{}万有之主天主将坐于其上；\newline{}其名可畏，我不敢\newline{}以软弱的人类言语呼出。}\par

\Person{欧里庇得斯}说：—\par

\Quote{祂不吝惜地赐予生命租期，\newline{}好让我们这些承租者受公正审判；\newline{}若有一凡人以为能隐藏\newline{}日常罪愆于天主之锐眼，\newline{}那是恶念；或许偶尔\newline{}他遇见悠闲的正义，她\newline{}便索要他为合法囚徒：\newline{}但你们中许多人草率地\newline{}犯下双重罪，说没有天主。\newline{}然而，啊！有天主；有天主。那么要留意\newline{}那作恶却亨通之人，或许能赎回\newline{}如此宝贵的时间，否则日后\newline{}必有相应的惩罚之报等候他。}\par

% structure: p0022 source=h2 target=section reason=relative-heading-rank; base=h2

\section{第四章　天主不愿祭献，只愿公义}

天主并不因恶人的奠酒与馨香而息怒，而是按公义报应各人，对此\Person{斐利蒙}又将为我作证：—\par

\Quote{若有人幻想，哦\Person{潘菲卢}，\newline{}通过献牛或山羊——不，那么，\newline{}凭朱庇特——或任何此类之物；\newline{}或奉献金银紫袍，\newline{}或象牙宝石雕像；\newline{}若他以为如此可平息天主之怒，\newline{}他错了，显示自己是愚人。\newline{}倒不如让他有益而良善，\newline{}不犯偷窃、淫行，\newline{}不为地上财富行凶杀人。\newline{}让他不觊觎他人之妻或子，\newline{}其华宅、广厦，\newline{}其少女、家生之奴，\newline{}其马匹、牲畜、牛群，\newline{}是的，连一根针也不贪，哦\Person{潘菲卢}，\newline{}因为天主就在你身边，看见你所行的一切。\newline{}祂常与恶人发怒，\newline{}却始终喜悦义人，\newline{}容许他收获劳作的果实，\newline{}享受汗水挣来的面包。\newline{}但你既为义人，务必偿还誓愿，\newline{}向施与者天主奉献礼物。\newline{}将你的装饰不放在外表，\newline{}而在内心的纯洁；\newline{}那时你听雷声，必不惧怕，\newline{}也不会在其声中逃遁，主人啊，\newline{}因为你自觉无恶行，\newline{}而天主就在你身边，看见你所行的一切。}\par

再者，\Person{柏拉图}在\WorkTitle{蒂迈欧篇}中说：\Quote{但若有人经过思考而真正进行严格的探究，他将对人性与神性之区分一无所知；因为天主将许多事物混而为一，[又能再次将一分解为许多事物，] 因祂赋有知识与大能；但无人能，也永不能行其中的任何一件。}\par

% structure: p0026 source=h2 target=section reason=relative-heading-rank; base=h2

\section{第五章　假神之虚荣僭妄}

至于那些以为能分享圣洁完全之名的人——有人凭虚妄传统接纳他们如神一般——\Person{米南德}在\WorkTitle{御者}中说：—\par

\Quote{若有一位神与老妇同行，\newline{}或偷偷潜入\newline{}折叠之门进入房宅，\newline{}他绝不能取悦我；不，唯有那一位\newline{}居家中、公义正直的天主，\newline{}为敬拜祂者赐予救恩。}\par

同一位\Person{米南德}在\WorkTitle{司祭}中说：—\par

\Quote{没有天主，啊妇人，能藉他人\newline{}拯救一人；若真有人\newline{}以铙钹之声随意迷惑神，\newline{}那么行此事者当然是更大的神。\newline{}但这些，\newline{}哦\Person{罗德}，不过是狡猾的计谋，\newline{}那些大胆、肆无忌惮的阴谋家\newline{}编造出来谋生，\newline{}成为时代的笑柄。}\par

同样，同一位\Person{米南德}陈述他对那些被接纳为神之人的意见，实则证明他们并非如此，他说：—\par

\Quote{是的，若我目睹此事，我便愿\newline{}我的灵魂返回我身。\newline{}因为告诉我，世上何处，哦\Person{格塔斯}，\newline{}可能找到公义的神？}\par

在\WorkTitle{寄存}中：—\par

\Quote{看来即便在众神之中，\newline{}也有不公的审判。}\par

悲剧家\Person{欧里庇得斯}在\WorkTitle{俄瑞斯忒斯}中说：—\par

\Quote{阿波罗藉其命令\newline{}导致弑母，却不知\newline{}何为诚实与正义。\newline{}我们侍奉众神，无论他们是谁；\newline{}但从大地的中央区域\newline{}你分明见阿波罗向凡人\newline{}发出神谕，凡他所说的我们都遵行。\newline{}我服从了他，当生我者\newline{}倒在我手下：他才是恶人。\newline{}那么杀他吧，因为犯罪的是他，不是我。\newline{}我能做什么？你难道不认为神\newline{}应免除我所担的罪责？}\par

同一位在\WorkTitle{希波吕托斯}中：—\par

\Quote{但在这些事上，众神并不公正判断。}\par

在\WorkTitle{伊翁}中：—\par

\Quote{但厄瑞克透斯的女儿\newline{}与我何干？因为那事\newline{}不属我辈。但当我前来\newline{}以金器行奠酒礼，我\newline{}将洒下露水，却仍须警告\newline{}阿波罗其行径；因他强娶\newline{}少女，秘密生子\newline{}后遗弃，任其死去。\newline{}故你并非如此；但你当趁时\newline{}常行美德，因为众神\newline{}必定惩罚恶人。\newline{}你为人类制定了律法，\newline{}自己却行不义，这如何合理？\newline{}但你既不在场，我将自由发言，\newline{}你须向人赔偿\newline{}强逼婚姻，你涅普顿、朱庇特，\newline{}统治天界的。你已使\newline{}神殿空虚，同时以不义回报。\newline{}虽你赞明智者为天上，\newline{}却以邪恶充满双手。\newline{}若人效仿众神之罪，\newline{}不可再称人为恶；\newline{}不，倒该称其师傅为恶。}\par

而在\WorkTitle{阿尔刻拉奥斯}中：—\par

\Quote{我儿，众神常使世人困惑。}\par

而在\WorkTitle{柏勒洛丰}中：—\par

\Quote{不行正义者，非神也。}\par

又同一剧中：—\par

\Quote{众神确在天上为王，人说；\newline{}但此言虚假，—且让那\newline{}如此说者，莫要愚蠢到\newline{}信古老传统，或留意\newline{}我言：当以无遮之眼\newline{}观此事于最清晰之光中。\newline{}我言之绝对权力，剥夺人的生命\newline{}与财产；背弃誓约；\newline{}亦不饶恕城市，反以残忍之手\newline{}无情劫掠摧毁。\newline{}然行此事者，反比\newline{}度温和虔诚生活者更为成功；\newline{}我知，那些敬畏神的小城，\newline{}屈从于更大不敬者的众多长枪；\newline{}是啊，我想\newline{}若有人闲荡懒惰，\newline{}不亲手劳作以求生计，\newline{}却向众神祈祷；他很快会知道\newline{}神是否会为他免除灾祸。}\par

而墨南德在\WorkTitle{迪菲卢斯}中：—\par

\Quote{因此我们当将赞美与尊荣归于\newline{}祂，万有之父，万有之主；\newline{}人类唯一的创造者与保存者，\newline{}祂以一切美物充满大地。}\par

同一作者在\WorkTitle{渔夫们}中：—\par

\Quote{因我认为那滋养我生命的\newline{}即是神；但那惯于满足\newline{}人之所需者—祂并不需要我们\newline{}重新供给，祂自身即是万有。}\par

同一作者在\WorkTitle{兄弟们}中：—\par

\Quote{神永远是正义者的\newline{}智慧：最智慧者如此认为。}\par

而在\WorkTitle{女笛手}中：—\par

\Quote{良好理性在万有中找到\newline{}敬拜的殿；因为心灵是什么，\newline{}岂非放置在我们内的神之声？}\par

而悲剧诗人在\WorkTitle{佛里克索斯}中：—\par

\Quote{若虔诚者与不虔者\newline{}\OriginalTerm{同命运}{share the same lot}，我们怎能认为公正，\newline{}如果至善的朱庇特，审判不公？}\par

在\WorkTitle{菲洛克忒忒斯}中：—\par

\Quote{你看见利益在神眼中\newline{}何等可敬；那些神龛\newline{}堆满黄金者何等受钦羡。\newline{}那么，既然像神一样是善的，\newline{}什么阻止你这样接受利益？}\par

在\WorkTitle{赫卡柏}中：—\par

\Quote{啊，朱庇特！无论你是谁，\newline{}除了名字，我们全然不知；}\par

以及—\par

\Quote{朱庇特，无论你确实是\newline{}伟大的必然，还是人的心灵，\newline{}我敬拜你！}\par

% structure: p0063 source=h2 target=section reason=relative-heading-rank; base=h2

\section{第六章　我们当承认唯一天主}

这里，便是德行的明证，也是爱慕智慧之心的明证：回归唯一者的共融，为得救而依附智慧，并按照人内所放置的自由意志选择更好的事物；而不认为那些具有人类情欲者是万有的主，因为他们甚至不会显明有与人同等的能力。因为在荷马笔下，得摩多科斯自称是无师自通的—\par

\Quote{神以灵感充满我}\par

—尽管他是一个凡人。阿斯克勒庇俄斯和阿波罗被半人马凯隆教导医术，—这可真是前所未闻的事，神竟被人教导。我何必说巴克斯，诗人说他发狂？或者赫拉克勒斯，说他是不幸的？何必说玛尔斯和维纳斯，通奸的领袖；借所有这些来确立我们已承担的证据？因为若有人无知，去模仿那些被说成是神圣的行为，他必被算为不洁之人，与生命和人性相异；若有人明知而为，他将有合理的借口逃避惩罚，表明模仿神明的大胆行为绝非罪过。但若有人责备这些行为，他必会除去它们众所周知的名称，而不会用华而不实的言辞掩盖。因此，必须接纳真实而不变的名号，不仅借我的言语宣告，也借那些将我们引入学问初阶者的话语宣告，为使我们不因在此现世生活中闲散度日，不仅如那些不知天上荣耀的人，而且显得忘恩负义，向审判者交账。\par

\SourceRecord{Translated by George Reith.；Ante-Nicene Fathers；Vol. 1.；Edited by Alexander Roberts, James Donaldson, and A. Cleveland Coxe.；Buffalo, NY: Christian Literature Publishing Co.,；1885.；<http://www.newadvent.org/fathers/0130.htm>.}
